%
Når du tråkker på bremsepedalen i en bil, presses et lite stempel inn i en sylinder med bremsevæske. Røret med bremsevæske fører videre til et stempel ved hvert av de fire hjulene. Stempelet ved pedalen har areal \SI{1.5}{\centi\metre\squared}, og hvert av stemplene ved hjulene har areal \SI{9.0}{\centi\metre\squared}. Kraften fra pedalen på det lille stempelet er \SI{200}{\newton}.
\begin{enumerate}[a)]
    \item Finn trykkøkningen i bremsevæsken.
    \item Finn kraften fra bremsevæsken på hvert av stemplene ved hjulene.
    \item Rørene fram til hjulene har ulik lengde. Hvorfor blir likevel kraften på stempelet like stor ved alle fire hjulene?
\end{enumerate}
%
